\chapter{Wavelets and Multiresolution}

In Chapter~10, the short-time Fourier transform let us see \emph{when} each frequency
occurs by chopping a signal into sliding windows. Now that Chapter~22 has given us
multirate filter banks, we can fix the spectrogram's central limitation: it made one window-length decision and
froze it. We picked a window length, and that single choice fixed our time
resolution and our frequency resolution \emph{everywhere}, for \emph{every}
frequency, for the whole signal. A short window let us pin down the timing of a
quick click but blurred the pitch of a slow hum. A long window resolved the hum
well but smeared the click across half a second. One window length meant one
compromise for the whole signal, even though different parts of a signal need
different resolutions.

Wavelets remove that compromise. The core idea
is simple: \emph{use short windows to look
at high frequencies and long windows to look at low frequencies}. Fast things
deserve a fast look; slow things deserve a patient one. This is
\textbf{multiresolution} -- analysing a signal at many scales at once, each scale
with the resolution appropriate to it. By the end of this chapter you will see
wavelets geometrically (as projections onto a family of stretched, shifted little
waves), you will work the Haar wavelet out by hand on real numbers, you will
understand the filter-bank machine that makes the discrete wavelet transform run
in linear time, and you will see why all of this is well suited to
catching the sharp \emph{onset} of a gesture on a wrist sensor -- a job the
Fourier transform is structurally bad at.

\chapterline

\section{The limits of a single window}
\label{ch24:tyranny}

Recall how the Short-Time Fourier Transform
(STFT) works. It slides a fixed window $w(t)$ of some length $T_w$ along the
signal and takes a Fourier transform inside each placement:
\begin{equation}
\label{ch24:stft}
\text{STFT}\{x\}(\tau, f) = \int_{-\infty}^{\infty} x(t)\, w(t-\tau)\,
e^{-\jj 2\pi f t}\,\dd t.
\end{equation}
The window has a width in time, call it $\Delta t$, and because of the
uncertainty principle its Fourier transform has a width in frequency
$\Delta f$, with
\begin{equation}
\label{ch24:uncertainty}
\Delta t \,\cdot\, \Delta f \;\ge\; \frac{1}{4\pi}.
\end{equation}
You cannot make both small. Worse, in the STFT the \emph{same} $\Delta t$ and the
same $\Delta f$ apply to the lowest frequency bin and the highest frequency bin
alike. The time-frequency plane gets tiled by identical rectangles, all the same
shape (Figure~\ref{ch24:tiling}, left).

\begin{intuitionbox}
Picture the time-frequency plane as a sheet of graph paper you are trying to
cover with tiles, where each tile is one ``measurement'' with some width in time
and some height in frequency. The uncertainty principle says every tile must have
at least a certain \emph{area}. The STFT forces every tile to be the same
\emph{shape} -- a fixed rectangle. Wavelets keep the area but vary the shape:
tall-and-thin tiles up high (good timing, coarse pitch), short-and-wide tiles
down low (good pitch, coarse timing). Same total cost, far better fit to how real
signals are organized.
\end{intuitionbox}

Why does the dyadic wavelet tiling suit real signals? Because of how their
information is spread across scales. A high-frequency event -- a tap, a
plosive in speech, the snap when your wrist flicks -- is necessarily brief; it is
made of fast wiggles and it is over quickly, so you want fine \emph{time}
resolution there and you can afford coarse frequency resolution (there is not
much room between ``high'' and ``very high'' that matters). A low-frequency
component -- the slow drift of your forearm, the fundamental of a bass note -- is
by its nature spread out in time, so fine \emph{frequency} resolution is what pays
off and coarse time resolution costs you nothing. Constant \emph{relative}
bandwidth, $\Delta f / f = \text{const}$, is the natural choice, and it is exactly
what musical octaves, the cochlea in your ear, and wavelets all use.

\begin{figure}[H]
\centering
\begin{tikzpicture}[scale=1.0, >=Stealth]
% --- STFT tiling (left) ---
\begin{scope}
\node[font=\small\bfseries, sourceblue!90!black] at (2.2, 4.8) {STFT: Uniform Time-Frequency Grid};
\draw[->, thick, gray!70] (0,0) -- (4.4,0) node[right, font=\footnotesize] {Time $t$};
\draw[->, thick, gray!70] (0,0) -- (0,4.2) node[left, font=\footnotesize] {Freq $f$};
\foreach \x in {0,1,2,3}{
  \foreach \y in {0,1,2,3}{
    \draw[sourceblue!80, fill=sourceblue!8, rounded corners=1pt] (\x+0.03,\y+0.03) rectangle (\x+0.97,\y+0.97);
  }}
\draw[<->, specred, thick] (1.1, 0.4) -- (1.9, 0.4);
\node[above, font=\tiny\bfseries, specred] at (1.5, 0.42) {$\Delta t$};
\draw[<->, specred, thick] (2.4, 0.1) -- (2.4, 0.9);
\node[right, font=\tiny\bfseries, specred] at (2.45, 0.5) {$\Delta f$};
\node[align=center, font=\footnotesize] at (2.1,-0.6) {Fixed window length $T_w$:\\$\Delta t$ and $\Delta f$ constant everywhere};
\end{scope}

% --- Wavelet tiling (right) ---
\begin{scope}[xshift=6.8cm]
\node[font=\small\bfseries, specred!90!black] at (2.2, 4.8) {Wavelets: Dyadic Multiresolution Tiling};
\draw[->, thick, gray!70] (0,0) -- (4.4,0) node[right, font=\footnotesize] {Time $t$};
\draw[->, thick, gray!70] (0,0) -- (0,4.2) node[left, font=\footnotesize] {Freq $f$};
% low freq band: 0-1, wide tiles
\foreach \x in {0,2}{
  \draw[specred!80, fill=specred!8, rounded corners=1pt] (\x+0.03,0.03) rectangle (\x+1.97,0.97);
}
% mid band 1-2: medium
\foreach \x in {0,1,2,3}{
  \draw[specred!80, fill=specred!12, rounded corners=1pt] (\x+0.03,1.03) rectangle (\x+0.97,1.97);
}
% high band 2-3: narrow
\foreach \x in {0,0.5,...,3.5}{
  \draw[specred!80, fill=specred!16, rounded corners=1pt] (\x+0.02,2.03) rectangle (\x+0.48,2.97);
}
% highest band 3-4: narrowest
\foreach \x in {0,0.25,...,3.75}{
  \draw[specred!80, fill=specred!22, rounded corners=1pt] (\x+0.01,3.03) rectangle (\x+0.24,3.97);
}
\node[align=center, font=\footnotesize] at (2.1,-0.6) {Scale-adaptive windows:\\Narrow $\Delta t$ at high $f$, wide $\Delta t$ at low $f$};
\end{scope}
\end{tikzpicture}
\caption{Tiling the time-frequency plane. The STFT (left) partitions the plane into identical rectangles, imposing a uniform resolution trade-off across all frequencies. The discrete wavelet transform (right) uses narrow, time-localized tiles at high frequencies (ideal for brief clicks, plosives, and gesture onsets) and broad, frequency-selective tiles at low frequencies (ideal for slow trends and musical pitch). Every tile satisfies the Heisenberg-Gabor uncertainty lower bound $\Delta t \cdot \Delta f \ge 1/(4\pi)$.}
\label{ch24:tiling}
\end{figure}

\begin{connectionbox}
The ``same area, different shape'' picture is the time--frequency uncertainty
principle from Chapter~7, which reduces algebraically to the Cauchy--Schwarz
inequality on inner products (Chapter~42). A signal and its
transform cannot both be arbitrarily concentrated, because concentration is an
inner-product statement and inner products obey Cauchy-Schwarz. Wavelets do not
beat the principle -- nothing does -- they just \emph{allocate} the unavoidable
blur intelligently.
\end{connectionbox}

\section{One little wave, stretched and shifted}
\label{ch24:one-wave}

The Fourier transform compares your signal against pure sinusoids $e^{\jj 2\pi f
t}$ -- infinite, eternal waves of a single frequency. That is why it has no sense
of time: its yardsticks have no location. A wavelet replaces the eternal sinusoid
with a \emph{little wave}: a function $\psi(t)$ that wiggles for a while and then
dies away to nothing (Figure~\ref{ch24:morlet}). It is local: it has a position.

From this single \emph{mother wavelet} $\psi$ we manufacture an entire family by
two operations: \textbf{scaling} (stretching or squishing it in time) and
\textbf{translation} (sliding it along):
\begin{equation}
\label{ch24:family}
\psi_{a,b}(t) \;=\; \frac{1}{\sqrt{a}}\,
\psi\!\left(\frac{t-b}{a}\right),
\qquad a > 0,\; b \in \R.
\end{equation}
Here $b$ is \emph{where} we look and $a$ is the \emph{scale} -- the stretch
factor. Large $a$ stretches $\psi$ out, so it wiggles slowly: a low-frequency
probe. Small $a$ squishes it, so it wiggles fast: a high-frequency probe. The
factor $1/\sqrt{a}$ out front is a normalization that keeps every member of the
family at the same energy, $\norm{\psi_{a,b}} = \norm{\psi}$, so that comparisons
across scales are fair.

\begin{happeningbox}
Consider what \eqref{ch24:family} does. As you make the probe wider (bigger
$a$) to chase lower frequencies, it automatically gets \emph{longer in time} --
it covers more of the signal. As you make it narrower (smaller $a$) to chase
higher frequencies, it automatically gets \emph{shorter in time}. The window
length is no longer a knob you set by hand; it is welded to the frequency you are
asking about. This coupling of window length to frequency is the multiresolution
idea; the rest of the chapter works out its consequences.
\end{happeningbox}

\begin{figure}[H]
\centering
\begin{tikzpicture}
\begin{axis}[
  width=12cm, height=4.4cm,
  xlabel={$t$}, ylabel={$\psi(t)$},
  xmin=-4, xmax=4, ymin=-1.1, ymax=1.1,
  axis lines=box, samples=180, domain=-4:4,
  title={A mother wavelet (real Morlet): a localized oscillation},
  title style={font=\small},
]
\addplot[sigblue, thick] {cos(deg(5*x))*exp(-x*x/2)};
\addplot[specred, dashed, thick] {exp(-x*x/2)};
\addplot[specred, dashed, thick] {-exp(-x*x/2)};
\end{axis}
\end{tikzpicture}
\caption{A mother wavelet $\psi(t) = \cos(5t)\,e^{-t^2/2}$: an oscillation
multiplied by a decaying envelope (dashed). Unlike a Fourier sinusoid it has a
definite location and finite reach. Stretching it (large $a$) makes it slower and
longer; squishing it (small $a$) makes it faster and shorter.}
\label{ch24:morlet}
\end{figure}

Two properties make a function a legal mother wavelet, and both are intuitive.
First, it must have \emph{zero average}, $\int \psi(t)\,\dd t = 0$ -- it wiggles
up as much as down, which is what ``wave'' means and which makes it blind to
constant offsets. Second, it must be localized (decay to zero), which is what
``little'' means. Together these are summarized by the \emph{admissibility
condition} $\int_0^\infty \abs{\Psi(f)}^2 / f \, \dd f < \infty$, where $\Psi$ is
the Fourier transform of $\psi$; the zero-mean requirement is exactly what keeps
this integral finite at $f=0$. In short, a wavelet is a zero-mean, localized
bump.

\section{The continuous wavelet transform}
\label{ch24:cwt}

Now we do the only thing we ever do in this book to measure ``how much of a
template is in a signal'': we take an inner product. The
\textbf{continuous wavelet transform} (CWT) of a signal $x(t)$ is its inner
product with every member of the wavelet family:
\begin{equation}
\label{ch24:cwt-defn}
W_x(a,b) \;=\; \inner{x}{\psi_{a,b}}
\;=\; \frac{1}{\sqrt{a}}\int_{-\infty}^{\infty}
x(t)\,\psi^{*}\!\left(\frac{t-b}{a}\right)\dd t.
\end{equation}
The output $W_x(a,b)$ is a function of two variables -- scale $a$ and position $b$
-- and it answers a concrete question: ``how strongly does the signal,
near time $b$, resemble a wiggle of scale $a$?'' Plot $\abs{W_x(a,b)}$ as an
image with $b$ across and $a$ up the side and you have a \textbf{scalogram}, the
wavelet cousin of the spectrogram.

\begin{whybox}
Why is this an inner product and not something fancier? Because the question ``how
much does $x$ look like the template $\psi_{a,b}$?'' is, geometrically, ``what is
the component of $x$ along the direction $\psi_{a,b}$?'' -- and the component of a
vector along a unit direction is exactly the inner product. The CWT is a
\emph{projection}, the same operation behind the Fourier transform, behind
least-squares, and inner products in Chapters~3 and~42. We have only
changed the set of directions we project onto: from eternal sinusoids to
localized, multi-scale wavelets. The machinery is identical; the yardsticks are
better suited to changing signals.
\end{whybox}

\begin{intuitionbox}
A useful mental animation: fix a scale $a$ and slide the wavelet across the signal
(vary $b$). You are doing a correlation -- a convolution with a flipped wavelet --
and you light up wherever the signal locally matches that wiggle. Now change $a$
and do it again with a stretched wavelet, lighting up where the signal has slower
structure. Stack all these slide-passes and you get the scalogram: fast structure
sits high (small $a$), slow structure sits low (large $a$), and each lights up at
the \emph{time} it occurs.
\end{intuitionbox}

\begin{applicationbox}
Scalograms work well on transients. A scratch on an old recording, a heartbeat's QRS
complex in an ECG, a crack in the acoustic emission of a stressed beam, the chirp
of a gravitational-wave inspiral (LIGO used wavelet-family methods in
its analysis pipelines): all are short, sharp, broadband events that a fixed-window
spectrogram either blurs in time or misses in frequency. The CWT catches them
because its high-frequency probes are short enough to pin the moment of the event
while its low-frequency probes still resolve the slow background.
\end{applicationbox}

\section{From continuous to discrete: the dyadic grid}
\label{ch24:dyadic}

The CWT computes $W_x(a,b)$ for \emph{every} real $a$ and $b$. That is hugely
redundant -- neighbouring scales and positions overlap enormously, so we are
storing the same information many times over. For computation and for compression
we want a discrete, non-redundant sample of the $(a,b)$ plane that still captures
everything. The standard choice is the \textbf{dyadic grid}: sample scale in
powers of two and sample position proportionally to scale.
\begin{equation}
\label{ch24:dyadic-grid}
a = 2^{j}, \qquad b = k\,2^{j}, \qquad j,k \in \Z.
\end{equation}
Each step in $j$ doubles the scale (one octave). At a coarser scale (larger $j$)
the wavelet is wider, so we slide it in correspondingly larger steps -- we do not
need to sample position finely when the probe itself is broad. The discretized
wavelets are
\begin{equation}
\label{ch24:dwt-basis}
\psi_{j,k}(t) = 2^{-j/2}\,\psi\!\left(2^{-j}t - k\right).
\end{equation}
The central mathematical discovery of multiresolution analysis is that,
for appropriately constructed wavelets $\psi$, this discrete, octave-spaced family forms
a complete, non-redundant \emph{orthonormal basis} for all finite-energy signals:
\begin{equation}
\label{ch24:orthonormal}
\inner{\psi_{j,k}}{\psi_{j',k'}} = \delta_{jj'}\,\delta_{kk'}.
\end{equation}
Every signal is a unique sum $x = \sum_{j,k} \inner{x}{\psi_{j,k}}\,\psi_{j,k}$,
with the wavelet coefficients $d_{j,k} = \inner{x}{\psi_{j,k}}$ playing exactly
the role the Fourier coefficients played -- except each one is local in both time
and scale. This is the \textbf{discrete wavelet transform} (DWT).

\begin{trapbox}
Do not confuse the DWT with ``the CWT sampled on a grid.'' Sampling the CWT
generally gives a redundant frame, not a basis. The DWT is a separate, sharper
object: a genuine orthonormal expansion that uses exactly $N$ coefficients to
represent $N$ samples, with a perfect inverse and an $O(N)$ algorithm. You get
that only for special, carefully engineered wavelets -- Haar, Daubechies, and
their kin -- not for an arbitrary localized bump.
\end{trapbox}

\section{The scaling function: where the leftovers go}
\label{ch24:scaling}

There is a loose end. Wavelets have zero mean -- they are deliberately blind to
the slow, near-constant trend of a signal. If we only ever project onto wavelets,
we can never reconstruct the average value or the lowest, slowest drift. We need a
partner that captures the leftovers: the \textbf{scaling function} $\varphi(t)$
(also called the \emph{father wavelet}), a localized bump with \emph{nonzero}
mean that represents the coarse, smoothed version of the signal.

The two functions are linked by two \emph{refinement equations}, which are
the core of the theory. They say that the scaling function at
one scale is a weighted combination of \emph{shifted copies of itself at the next
finer scale}:
\begin{align}
\label{ch24:refine-phi}
\varphi(t) &= \sqrt{2}\,\sum_{n} h[n]\,\varphi(2t-n), \\
\label{ch24:refine-psi}
\psi(t) &= \sqrt{2}\,\sum_{n} g[n]\,\varphi(2t-n).
\end{align}
The coefficients $h[n]$ and $g[n]$ are short sequences of numbers -- a low-pass
filter $h$ and a high-pass filter $g$ -- and \emph{they are all you
need}. Give me the filter taps $h[n]$ (with $g[n]$ derived from them by a simple
alternating-sign-and-reverse rule, $g[n] = (-1)^n h[1-n]$) and I can generate the
scaling function, the wavelet, and the entire transform. The continuous functions
$\varphi$ and $\psi$ are, in a real sense, just elaborate ways of drawing a
picture of a handful of filter taps.

\begin{happeningbox}
The refinement equation \eqref{ch24:refine-phi} expresses self-similarity as an
equation: ``I am built out of smaller copies of myself.'' That is why wavelets are
deeply related to fractals, and why a smooth-looking Daubechies scaling function
is actually generated by an infinitely-recursive subdivision of a few numbers. In
practice this matters a great deal: because the whole transform is encoded in short
filters, we never have to handle the continuous functions at all on a computer. We
just filter and downsample. The fact that continuous multiresolution
reduces to discrete filtering is what turned wavelets from a mathematical curiosity
into a practical everyday tool.
\end{happeningbox}

\section{Haar: the wavelet you can do by hand}
\label{ch24:haar}

The oldest and simplest wavelet, the
\textbf{Haar wavelet} (Alfred Haar, 1909 -- two decades before the word
``wavelet'' existed), is one you can compute with on paper. Its scaling function
is a single unit box, and its wavelet is a box that goes up then down:
\begin{equation}
\label{ch24:haar-defn}
\varphi(t) =
\begin{cases} 1 & 0 \le t < 1\\ 0 & \text{else}\end{cases}
\qquad\qquad
\psi(t) =
\begin{cases} +1 & 0 \le t < \tfrac12\\ -1 & \tfrac12 \le t < 1\\ 0 &
\text{else}.\end{cases}
\end{equation}

\begin{figure}[H]
\centering
\begin{tikzpicture}
% scaling function
\begin{scope}
\node[font=\small\bfseries, color=sigblue] at (1.0, 2.1) {Scaling Function $\varphi(t)$};
\draw[->] (-0.4,0)--(2.2,0) node[right]{\small $t$};
\draw[->] (0,-1.3)--(0,1.7) node[left]{\small $\varphi$};
\draw[sigblue,very thick] (0,1)--(1,1);
\draw[sigblue,very thick,dashed] (0,0)--(0,1);
\draw[sigblue,very thick,dashed] (1,0)--(1,1);
\draw (1, 0.06) -- (1, -0.06);
\node[below=2pt] at (1, 0) {\small $1$};
\node[left] at (0,1) {\small $1$};
\node[font=\footnotesize, text=gray!80!black] at (1.0,-1.5) {Local average};
\end{scope}
% wavelet
\begin{scope}[xshift=5.2cm]
\node[font=\small\bfseries, color=specred] at (1.0, 2.1) {Haar Wavelet $\psi(t)$};
\draw[->] (-0.4,0)--(2.2,0) node[right]{\small $t$};
\draw[->] (0,-1.3)--(0,1.7) node[left]{\small $\psi$};
\draw[specred,very thick] (0,1)--(0.5,1);
\draw[specred,very thick] (0.5,-1)--(1,-1);
\draw[specred,very thick,dashed] (0,0)--(0,1);
\draw[specred,very thick,dashed] (0.5,-1)--(0.5,1);
\draw[specred,very thick,dashed] (1,-1)--(1,0);
\node[above left=1pt and 0pt, font=\footnotesize] at (0.5, 0) {$\tfrac{1}{2}$};
\draw (1, 0.06) -- (1, -0.06);
\node[above=2pt] at (1, 0) {\small $1$};
\node[left] at (0,1) {\small $1$};
\node[left] at (0,-1) {\small $-1$};
\node[font=\footnotesize, text=gray!80!black] at (1.0,-1.5) {Local difference};
\end{scope}
\end{tikzpicture}
\caption{The Haar scaling function (a box, left) and the Haar wavelet (up-then-down,
right). The box captures a local average; the up-then-down captures a local
difference. Multiresolution analysis is built from averages and differences.}
\label{ch24:haar-fig}
\end{figure}

What do these compute? Take a signal and a single Haar scaling function over an
interval: the inner product is (up to a constant) the \emph{average} of the signal
on that interval. Take the Haar wavelet: the inner product is the \emph{difference}
between the average of the first half and the average of the second half -- a
local \emph{detail}. The Haar transform computes only these
averages and differences.

The Haar filter taps are correspondingly the simplest possible. From
\eqref{ch24:refine-phi}, the box $\varphi(t)$ equals the sum of two
half-width boxes, $\varphi(t) = \varphi(2t) + \varphi(2t-1)$, which matched
against \eqref{ch24:refine-phi} gives
\begin{equation}
\label{ch24:haar-filters}
h = \left[\tfrac{1}{\sqrt2},\ \tfrac{1}{\sqrt2}\right],
\qquad
g = \left[\tfrac{1}{\sqrt2},\ -\tfrac{1}{\sqrt2}\right].
\end{equation}
The low-pass filter $h$ is ``add the pair and rescale'' (an average); the
high-pass filter $g$ is ``subtract the pair and rescale'' (a difference). The
$1/\sqrt2$ keeps energy constant.

\begin{example}[One step of the Haar DWT by hand]
\label{ch24:haar-example}
Take the eight-sample signal
\[
x = [\,4,\ 6,\ 10,\ 12,\ 10,\ 6,\ 4,\ 2\,].
\]
Process it in adjacent pairs. For each pair $(x_{2i}, x_{2i+1})$ compute a scaled
average $s_i = (x_{2i}+x_{2i+1})/\sqrt2$ and a scaled difference
$d_i = (x_{2i}-x_{2i+1})/\sqrt2$.

\textbf{Averages} (the coarse approximation):
\[
s = \tfrac{1}{\sqrt2}[\,4{+}6,\ 10{+}12,\ 10{+}6,\ 4{+}2\,]
  = \tfrac{1}{\sqrt2}[\,10,\ 22,\ 16,\ 6\,]
  \approx [\,7.07,\ 15.56,\ 11.31,\ 4.24\,].
\]
\textbf{Differences} (the detail coefficients):
\[
d = \tfrac{1}{\sqrt2}[\,4{-}6,\ 10{-}12,\ 10{-}6,\ 4{-}2\,]
  = \tfrac{1}{\sqrt2}[\,-2,\ -2,\ 4,\ 2\,]
  \approx [\,-1.41,\ -1.41,\ 2.83,\ 1.41\,].
\]
We have turned eight numbers into four averages and four details -- still eight
numbers, no information lost (the transform is invertible: recover each pair from
$x_{2i} = (s_i+d_i)/\sqrt2$, $x_{2i+1} = (s_i-d_i)/\sqrt2$). Notice the details are
\emph{small} where the signal is smooth and \emph{large} where it changes fast.
Those small values are what make compression possible.
\end{example}

\begin{example}[The full multiresolution pyramid]
\label{ch24:haar-pyramid}
The DWT does not stop after one step -- we recurse on the
\emph{averages}, leaving the details aside. Continuing Example~\ref{ch24:haar-example}
with $s \approx [7.07, 15.56, 11.31, 4.24]$:

Level 2 averages:
\[
\tfrac{1}{\sqrt2}[7.07{+}15.56,\ 11.31{+}4.24] \approx [16.0,\ 11.0],
\qquad
\tfrac{1}{\sqrt2}[7.07{-}15.56,\ 11.31{-}4.24] \approx [-6.0,\ 5.0]
\]
(the second pair being the level-2 details).

Level 3: average $\tfrac{1}{\sqrt2}[16.0{+}11.0] \approx 19.09$; detail
$\tfrac{1}{\sqrt2}[16.0{-}11.0] \approx 3.54$.

The final transform stacks one coarsest average with all the details from every
level:
\[
[\underbrace{19.09}_{\text{coarse}},\ \underbrace{3.54}_{\text{lvl 3}},\
\underbrace{-6.0,\ 5.0}_{\text{lvl 2}},\
\underbrace{-1.41,\ -1.41,\ 2.83,\ 1.41}_{\text{lvl 1}}\,].
\]
Eight inputs, eight outputs, organized from coarsest to finest. The single coarse
number is (a scaled version of) the overall average; everything else is detail at
ever-finer time scales. \emph{This} is multiresolution: the signal decomposed into
a coarse sketch plus a hierarchy of corrections.
\end{example}

\begin{recoverybox}
\textbf{Pause and think}: Why does the transform produce exactly one coarsest average coefficient, but four level-1 detail coefficients? Each time you halve the scale, the wavelet shrinks by half, so you need twice as many translated copies to span the signal. Level 1 details check differences across adjacent 2-sample blocks (4 coefficients); level 2 checks differences across 4-sample blocks (2 coefficients); level 3 checks differences across the whole 8-sample signal (1 coefficient). Summing them: $1 + 1 + 2 + 4 = 8$ numbers, preserving the exact dimension of the input vector space.
\end{recoverybox}

\begin{figure}[H]
\centering
\begin{tikzpicture}[font=\small]
\node[draw,fill=sigblue!10,minimum width=6.4cm,minimum height=0.7cm] (s0) at (0,3) {$x$: 8 samples};
\node[draw,fill=sigblue!14,minimum width=3.2cm,minimum height=0.7cm] (s1) at (-1.6,2) {avg (4)};
\node[draw,fill=specred!12,minimum width=3.0cm,minimum height=0.7cm] (d1) at (1.7,2) {detail $d_1$ (4)};
\node[draw,fill=sigblue!18,minimum width=1.6cm,minimum height=0.7cm] (s2) at (-2.4,1) {avg (2)};
\node[draw,fill=specred!12,minimum width=1.5cm,minimum height=0.7cm] (d2) at (-0.85,1) {$d_2$ (2)};
\node[draw,fill=sigblue!22,minimum width=0.9cm,minimum height=0.7cm] (s3) at (-2.85,0) {a};
\node[draw,fill=specred!12,minimum width=0.9cm,minimum height=0.7cm] (d3) at (-1.85,0) {$d_3$};
\draw[->] (s0) -- (s1); \draw[->] (s0) -- (d1);
\draw[->] (s1) -- (s2); \draw[->] (s1) -- (d2);
\draw[->] (s2) -- (s3); \draw[->] (s2) -- (d3);
\node[align=left,text width=4cm] at (3.7,0.5){Final coefficients:\\ one coarse $a$ +\\ details $d_3,d_2,d_1$};
\end{tikzpicture}
\caption{The multiresolution pyramid. Each level low-passes (averages) and
high-passes (differences) the previous approximation, then keeps the details and
recurses only on the average. The details accumulate; the approximation shrinks by
half each level. This cascade is the discrete wavelet transform.}
\label{ch24:pyramid}
\end{figure}

\section{The filter-bank view: it is just filter, downsample, repeat}
\label{ch24:filterbank}

Examples~\ref{ch24:haar-example} and \ref{ch24:haar-pyramid} already contain the
algorithm. Each level does three things:
\begin{enumerate}
\item convolve the current approximation with the low-pass filter $h$ (averaging),
\item convolve it with the high-pass filter $g$ (differencing),
\item \textbf{downsample} both outputs by 2 (keep every other sample -- we only
needed one average and one difference per pair).
\end{enumerate}
Keep the high-pass output as that level's detail; feed the low-pass output into
the next level. This is \textbf{Mallat's fast wavelet transform} (Stephane Mallat
and Yves Meyer, 1989), and it is what made wavelets practical.

\begin{figure}[H]
\centering
\begin{tikzpicture}[font=\small, >=Stealth]
\tikzset{
  flt/.style={draw=blue!70!black, fill=sourceblue!10, rounded corners=2pt, minimum height=8mm, minimum width=13mm, align=center, thick},
  hflt/.style={draw=specred!70!black, fill=specred!10, rounded corners=2pt, minimum height=8mm, minimum width=13mm, align=center, thick},
  ds/.style={draw=gray!70!black, fill=gray!10, circle, inner sep=1.5pt, thick, font=\footnotesize}
}

% Stage 1 boundary background
\draw[draw=gray!30, fill=gray!3, rounded corners=4pt, dashed] (0.8, -2.6) rectangle (5.3, 1.8);
\node[font=\footnotesize\scshape, text=gray!70] at (3.0, 1.5) {Octave 1 Decomposition};

% Stage 2 boundary background
\draw[draw=gray!30, fill=gray!3, rounded corners=4pt, dashed] (5.8, -2.6) rectangle (10.3, 1.8);
\node[font=\footnotesize\scshape, text=gray!70] at (8.0, 1.5) {Octave 2 Decomposition};

% Input
\node (x) at (0,0) [font=\normalsize\bfseries] {$x[n]$};

% Level 1
\node[flt] (h1) at (2.2, 0.7) {Lowpass\\$h[n]$};
\node[hflt] (g1) at (2.2, -1.2) {Highpass\\$g[n]$};
\node[ds] (dh1) at (4.2, 0.7) {$\downarrow 2$};
\node[ds] (dg1) at (4.2, -1.2) {$\downarrow 2$};

\draw[->, thick] (x) -- ++(1.0, 0) |- (h1.west);
\draw[->, thick] (1.0, 0) |- (g1.west);
\draw[->, thick] (h1) -- (dh1);
\draw[->, thick] (g1) -- (dg1);

\node (d1) at (4.2, -2.1) [below, font=\scriptsize\bfseries, specred] {$d_1[n]$ (Detail 1)};
\draw[->, thick, specred] (dg1) -- (4.2, -2.0);

% Level 2
\node[flt] (h2) at (7.2, 0.7) {Lowpass\\$h[n]$};
\node[hflt] (g2) at (7.2, -1.2) {Highpass\\$g[n]$};
\node[ds] (dh2) at (9.2, 0.7) {$\downarrow 2$};
\node[ds] (dg2) at (9.2, -1.2) {$\downarrow 2$};

\draw[->, thick] (dh1) -- (h2) node[midway, above, font=\tiny] {$a_1[n]$};
\draw[->, thick] (5.0, 0.7) |- (g2.west);
\draw[->, thick] (h2) -- (dh2);
\draw[->, thick] (g2) -- (dg2);

\node (d2) at (9.2, -2.1) [below, font=\scriptsize\bfseries, specred] {$d_2[n]$ (Detail 2)};
\draw[->, thick, specred] (dg2) -- (9.2, -2.0);

\node (a2) at (10.5, 0.7) [right, font=\scriptsize\bfseries, sourceblue] {$a_2[n]$ (Approx)};
\draw[->, thick, sourceblue] (dh2) -- (a2);

\end{tikzpicture}
\caption{Mallat's dyadic filter bank tree (analysis stage). At each scale octave, the signal approximation is split by lowpass ($h[n]$) and highpass ($g[n]$) quadrature mirror filters, then decimated by 2. The detail stream $d_j[n]$ is preserved as the wavelet coefficients for that octave, while the downsampled approximation $a_j[n]$ cascades into the next coarser stage.}
\label{ch24:filterbank-fig}
\end{figure}

\begin{whybox}
Why does downsampling not throw information away? After low-pass filtering, the
approximation occupies only the bottom half of the frequency band; by the sampling
theorem it is fully described at half the sample rate, so keeping every other
sample is lossless. Likewise the detail occupies the top half and survives its own
downsampling. The two half-rate streams together carry exactly the information of
the full-rate input -- this is the \emph{perfect reconstruction} property, and it
is why a careful filter-bank design (the filters must satisfy the
\emph{quadrature mirror} conditions -- meaning the high-pass filter is an
alternating-sign, time-reversed copy of the low-pass filter designed so their
aliasing terms cancel out exactly when recombined) loses nothing.
\end{whybox}

The cost is low. Level 1 filters $N$ samples; level 2 filters $N/2$; level
3 filters $N/4$; and so on. The total work is $N + N/2 + N/4 + \cdots < 2N$, so
the entire DWT is \textbf{$O(N)$} -- even faster than the FFT's $O(N\log N)$. A
multi-scale, time-localized analysis that runs in linear time is a huge practical
advantage on a small microcontroller.

\begin{historybox}
The ideas came from several fields over many decades. Alfred Haar constructed his
orthonormal step-function basis in 1909 to explore function spaces in pure mathematics.
In the early 1980s, the French exploration geophysicist Jean Morlet, analyzing reflected
acoustic pulses from underground seismic strata that the fixed-window STFT blurred,
introduced scaled-and-shifted probe wavelets of constant shape; he collaborated with
mathematical physicist Alex Grossmann to formulate the continuous wavelet transform.
Pure mathematicians initially doubted that smooth, orthogonal, compactly supported
wavelets could exist, until Yves Meyer discovered smooth orthonormal wavelets in 1985.
In 1988, Ingrid Daubechies solved the open algebraic problem by constructing families of
compactly supported, orthogonal wavelets with prescribed smoothness, while St\'ephane
Mallat connected these bases to the discrete multirate filter banks already familiar
in subband speech coding (1989), producing the $O(N)$ discrete pyramid algorithm.
This work led to standards like JPEG~2000 and Donoho--Johnstone wavelet shrinkage
denoising. An active frontier remains the development of directional, anisotropic
multiscale transforms (such as shearlets and curvelets) that represent curved edges in
higher-dimensional signals without the cross-term artifacts of separable 1D tensor products.
\end{historybox}

\begin{labbox}
PyWavelets makes the whole pyramid one call. Try it on a signal with a transient:
\begin{lstlisting}[language=Python]
import numpy as np, pywt
# smooth ramp with a sudden spike at n=300 (a "gesture onset")
n = np.arange(512)
x = 0.5*np.sin(2*np.pi*3*n/512)
x[300:305] += 4.0                       # the transient
coeffs = pywt.wavedec(x, 'db4', level=5)   # multilevel DWT
for j, c in enumerate(coeffs):
    print(f"band {j}: {len(c):4d} coeffs, max |d| = {np.abs(c).max():.2f}")
# the FINEST detail band lights up sharply right at the spike -- and
# its index pinpoints WHEN. The smooth sine barely touches the details.
xr = pywt.waverec(coeffs, 'db4')        # perfect reconstruction
print("reconstruction error:", np.max(np.abs(x - xr)))
\end{lstlisting}
Swap \texttt{'db4'} for \texttt{'haar'} to see the difference: Haar reacts
instantly but blockily; Daubechies-4 is smoother and localizes a clean spike with
fewer ringing artifacts.
\end{labbox}

\section{Wavelets for compression and denoising}
\label{ch24:compression}

Look again at the transformed signal in Example~\ref{ch24:haar-pyramid}: most of
the detail coefficients are small. This is typical of real
signals -- images, audio, sensor traces -- because real signals are mostly smooth
with occasional sharp features. The DWT \emph{concentrates the energy} into a few
large coefficients (the coarse trend and the details at genuine edges) and leaves
a sea of near-zero coefficients elsewhere. This is called \textbf{sparsity}, and
it is the foundation of both compression and denoising.

\textbf{Compression}: keep the largest coefficients, throw away (zero out) the
tiny ones, and store only the survivors. Reconstruct from those; the error is
tiny because you discarded only small numbers. This is precisely how
\textbf{JPEG~2000} compresses images -- a wavelet transform, then quantize and
keep the significant coefficients -- and it visibly beats old block-DCT JPEG at
high compression, with graceful blur instead of ugly blocky artifacts (the FBI
adopted a wavelet scheme for its fingerprint database for exactly this reason).

\textbf{Denoising}: white noise spreads its energy \emph{evenly} across all
wavelet coefficients, so it shows up as a uniform low-level fuzz, while the real
signal lives in a few large coefficients. \emph{Thresholding} -- zeroing every
coefficient below some level, then inverting the transform -- erases the fuzz and
keeps the signal. Donoho and Johnstone's \emph{wavelet shrinkage} (early 1990s)
made this rigorous and near-optimal. Unlike a linear low-pass filter,
wavelet thresholding suppresses background noise while keeping sharp edges intact,
because edge transitions survive as isolated large detail coefficients that comfortably exceed
the threshold.

\begin{figure}[H]
\centering
\begin{tikzpicture}
\begin{axis}[
  width=12cm, height=4.6cm,
  xlabel={coefficient index (coarse $\to$ fine)},
  ylabel={$\abs{\text{coeff}}$},
  xmin=-0.5, xmax=15.5, ymin=0, ymax=10,
  axis lines=box, ybar, bar width=6pt,
  xtick={0,4,8,12}, enlarge x limits=0.03,
  title={Energy compaction: a few large coefficients, the rest near zero},
  title style={font=\small},
]
\addplot[draw=sigblue,fill=fillsig] coordinates {
 (0,9.2)(1,4.1)(2,0.3)(3,0.2)(4,2.8)(5,0.15)(6,0.1)(7,3.4)
 (8,0.2)(9,0.1)(10,0.05)(11,0.3)(12,0.1)(13,0.05)(14,2.1)(15,0.1)};
\end{axis}
\end{tikzpicture}
\caption{A typical wavelet coefficient magnitude profile for a piecewise-smooth
signal. A handful of coefficients (the coarse trend, and details sitting on real
edges) carry almost all the energy; the rest are nearly zero. Keeping the big ones
compresses; thresholding away the small ones denoises.}
\label{ch24:compaction}
\end{figure}

\begin{connectionbox}
Sparsity connects wavelets directly to statistical estimation (Chapter~20). ``Most
coefficients are tiny'' is a statement about a heavy-tailed coefficient
distribution; choosing a threshold is a hypothesis test (signal vs.\ noise) per
coefficient; and the optimal threshold $\sigma\sqrt{2\ln N}$ for $N$ samples of
noise of standard deviation $\sigma$ comes straight from the expected maximum of
$N$ Gaussian draws. Denoising is therefore a statistical estimation problem.
\end{connectionbox}

\section{Transients, onsets, and the wrist controller}
\label{ch24:onsets}

This is why wavelets matter for the wrist controller. The gestures that drive our
instrument controller are defined by their \emph{transitions}: the instant a flick
begins, the snap of a tap, the moment a hold is released. These are \emph{onsets}
-- short, sharp, broadband events sitting on top of slow background motion. They
are precisely the signals the Fourier transform handles worst and the wavelet
transform handles best.

\begin{figure}[H]
\centering
\begin{tikzpicture}
\begin{axis}[
  width=12cm, height=3.4cm,
  xlabel={}, ylabel={accel mag.},
  xmin=0, xmax=10, ymin=-1.4, ymax=1.6,
  axis lines=left, samples=160, domain=0:10,
  xtick=\empty,
  title={Top: wrist signal -- slow sway with a flick at $t\approx 6$},
  title style={font=\small},
]
\addplot[sigblue, thick] {0.5*sin(deg(0.9*x)) + 0.9*sin(deg(18*(x-6.0)))*exp(-6*(x-6.0)^2)};
\end{axis}
\begin{axis}[
  width=12cm, height=3.0cm,
  yshift=-3.0cm,
  xlabel={$t$}, ylabel={finest detail},
  xmin=0, xmax=10, ymin=-0.1, ymax=1.0,
  axis lines=left, samples=160, domain=0:10,
  title={Bottom: finest DWT detail band -- near zero, then spikes at the onset},
  title style={font=\small},
]
\addplot[specred, thick] {0.02 + 0.85*exp(-30*(x-6.0)^2)};
\end{axis}
\end{tikzpicture}
\caption{Why the DWT is a natural onset detector. The slow forearm sway lives in
the coarse bands and barely touches the finest detail band, which stays near zero
(bottom) until the flick begins -- then it spikes, and the position of the spike
pinpoints \emph{when}. Thresholding this band gives a crisp gesture trigger.}
\label{ch24:onset-fig}
\end{figure}

\begin{goalbox}
On a wrist-worn IMU, a deliberate flick produces a brief, high-frequency burst in
the accelerometer and gyro -- riding on the slow, low-frequency sway of the whole
forearm. Run a few-level DWT on the incoming stream and watch the \emph{finest
detail band}: it stays near zero during smooth motion and \emph{spikes} the
instant a flick starts, with the coefficient's position telling you \emph{when}
down to a handful of samples (Figure~\ref{ch24:onset-fig}). That is an onset detector with built-in time
resolution, far crisper than waiting for an STFT frame to fill. A simple,
effective gesture trigger: threshold the energy in the finest one or two detail
bands; when it exceeds a baseline, fire a MIDI note-on. The coarse approximation
band, meanwhile, gives you the slow posture/tilt for free in the same single
transform -- onset \emph{and} pose from one $O(N)$ pass.
\end{goalbox}

\begin{buildbox}
A Haar (or short Daubechies) DWT is cheap enough to run live on an Arduino-class
MCU at IMU rates (say 200\,Hz). You do not even need the textbook pyramid: a
\emph{lifting} implementation of Haar is just averages and differences in place,
no extra buffers. Sketch of the onset stage:
\begin{lstlisting}[language=C]
// one Haar lifting step over a small ring buffer of accel magnitude
// (low-pass = average, high-pass = difference), in place
for (int i = 0; i < N; i += 2) {
    float avg  = 0.70710678f * (buf[i] + buf[i+1]);
    float diff = 0.70710678f * (buf[i] - buf[i+1]);
    approx[i/2] = avg;       // feeds next level (slow posture)
    detail[i/2] = diff;      // finest detail = onset energy
}
// onset trigger: short-term energy of finest detail vs. baseline
float e = 0;
for (int k = 0; k < N/2; k++) e += detail[k]*detail[k];
if (e > GATE * baseline) midi_note_on();   // fire the gesture
baseline = 0.98f*baseline + 0.02f*e;        // slow adaptive floor
\end{lstlisting}
Tune \texttt{GATE} on real captures of your own flicks; the adaptive baseline
keeps it working as the resting noise floor drifts.
\end{buildbox}

\begin{trapbox}
A reality check. A wrist IMU detecting a sharp onset
tells you that \emph{something abrupt happened} and roughly its intensity and
timing. It does \emph{not} tell you which finger moved or by how much. The DWT is
a good \emph{onset and energy} detector; it does not conjure information the
sensor never captured. Distinguishing individual fingers needs the muscle-level
signals of EMG or force myography on the forearm; continuous finger \emph{pose}
needs sensors on the fingers themselves (rings or gloves). Wavelets sharpen
\emph{when} and \emph{how hard}; they cannot upgrade a wrist IMU into a finger
tracker.
\end{trapbox}

\begin{aestheticbox}
Compare these time--frequency tools. The Fourier transform resolves
frequency globally while giving up time. The STFT recovers time at a single fixed window scale.
Wavelets let the window scale vary with frequency -- short windows for fast transients,
long windows for slow oscillations -- matching logarithmic octave spacing in music and
the mechanical frequency map of the human cochlea. Stretching and shifting a single prototype
wavelet yields an orthonormal basis, an $O(N)$ filter-bank cascade, subband image compression,
shrinkage denoising, and a sharp transient detector from one construction.
\end{aestheticbox}

Having seen how wavelets decompose signals across octaves of scale, we turn in the next
chapter to the Hilbert transform and the analytic signal. While wavelets adapt time-frequency
resolution, the Hilbert transform shifts the phase of every positive frequency by
$-90^\circ$ without changing its amplitude, giving a rigorous meaning to instantaneous
amplitude, phase, and frequency.

\section{Exercises}
\label{ch24:exercises}

\begin{exercise}
\label{ch24:ex-tiling}
Sketch the time-frequency tiling for an STFT with a 64-sample window and for a
4-level dyadic wavelet transform of the same 256-sample signal. For each, state
the time resolution and frequency resolution in the lowest and highest analyzed
bands. Explain in one sentence why the wavelet tiling is better suited to a signal
that is a slow hum with an occasional sharp click.
\end{exercise}

\begin{exercise}
\label{ch24:ex-haar-byhand}
Compute the full Haar DWT (all levels) of $x = [2, 4, 4, 6, 8, 8, 6, 4]$ by hand,
using the scaled average/difference rule from Example~\ref{ch24:haar-example}.
Report the final coefficient vector (coarsest average followed by details from
coarse to fine). Verify it has the same number of entries as $x$.
\end{exercise}

\begin{exercise}
\label{ch24:ex-invert}
Take your level-1 averages and details from the previous exercise and reconstruct
the original eight samples using $x_{2i} = (s_i+d_i)/\sqrt2$ and
$x_{2i+1} = (s_i-d_i)/\sqrt2$. Confirm you recover $x$ exactly. What does this
demonstrate about the DWT as a linear map?
\end{exercise}

\begin{exercise}
\label{ch24:ex-energy}
Show numerically (using your answer to Exercise~\ref{ch24:ex-haar-byhand}) that
the sum of squares of the wavelet coefficients equals the sum of squares of the
original samples. Which property of the Haar transform guarantees this, and what
is the name of the analogous Fourier result?
\end{exercise}

\begin{exercise}
\label{ch24:ex-zeromean}
The Haar wavelet $\psi$ has integral zero, but the scaling function $\varphi$ does
not. Compute both integrals from \eqref{ch24:haar-defn}. Explain why a signal's
overall average is carried entirely by the coarsest scaling coefficient and never
by any wavelet coefficient.
\end{exercise}

\begin{exercise}
\label{ch24:ex-filters}
Verify that the Haar filters in \eqref{ch24:haar-filters} satisfy
$\sum_n h[n]^2 = 1$ (unit energy) and $\sum_n h[n]\,g[n] = 0$ (orthogonality of
the low- and high-pass channels). Why are these two conditions exactly what a
perfect-reconstruction, orthonormal filter bank needs?
\end{exercise}

\begin{exercise}
\label{ch24:ex-cost}
A signal has $N = 2^{20}$ samples. Counting one multiply-add per filter tap, give
the approximate total operation count for a full Haar DWT and compare it to an
FFT of the same signal ($\approx N\log_2 N$). By what factor is the DWT cheaper?
Why does the geometric series $N + N/2 + N/4 + \cdots$ matter to your answer?
\end{exercise}

\begin{exercise}
\label{ch24:ex-denoise}
You add white Gaussian noise of standard deviation $\sigma = 0.1$ to a signal of
$N = 1024$ samples and take its DWT. Using the universal threshold
$\lambda = \sigma\sqrt{2\ln N}$, compute $\lambda$. Describe the denoising
procedure (threshold then invert) in three steps, and explain why this removes
noise without blurring a sharp edge in the signal, unlike a low-pass filter.
\end{exercise}

\begin{exercise}
\label{ch24:ex-pywt}
Using PyWavelets, generate a 3\,Hz sine of length 512 with a 5-sample spike added
at index 200. Run \texttt{pywt.wavedec} with \texttt{'db2'} at 4 levels. Plot each
detail band stacked. Identify which band and which coefficient index pinpoint the
spike, and confirm the spike's location matches index 200 after accounting for the
downsampling at each level.
\end{exercise}

\begin{exercise}
\label{ch24:ex-onset}
Design (on paper) a wrist-IMU gesture-onset detector based on the finest Haar
detail band, following the C onset sketch in Section~\ref{ch24:onsets}. Specify: the input signal you would
feed it (which IMU channel(s) and any preprocessing), how you would set the
threshold from a calibration recording, and one false-trigger failure mode (e.g.\
bumping the table). State plainly one thing this detector \emph{cannot} tell you
about the hand, and what sensor you would add to get it.
\end{exercise}

\begin{exercise}
\label{ch24:ex-relative}
The wavelet transform has constant \emph{relative} bandwidth ($\Delta f / f$
fixed), while the STFT has constant \emph{absolute} bandwidth ($\Delta f$ fixed).
Compute the absolute bandwidth of a wavelet band centered at 10\,Hz and one
centered at 160\,Hz if $\Delta f / f = 0.3$. Explain why this matches the way
musical pitch and human hearing are organized, and why an STFT with a single
window cannot reproduce it.
\end{exercise}
